\section{Additional Results}
\label{app:additional}

\subsection{Full hypothesis ledger}

Table~\ref{tab:ledger} in the main text lists all nine hypotheses. Two procedural notes belong
with it.

First, in two places our scoring code and our written decision rule diverged. In H8 the
pre-registration assigned \textsc{not supported} only when the first prediction failed and
\textsc{partial} when the first two held, leaving the case ``first holds, second fails''
unspecified; the code defaulted to \textsc{not supported} and we report that rather than
choosing a label after seeing the outcome. In H9 the pre-registration states that the verdict
is evaluated only over configurations passing the validity gate, which would make two
predictions uninformative, while the code scored one of them as a failure. Both paths reach the
same verdict, so nothing changes, but we record the discrepancy.

Second, the ensemble-size ablation (seed 0) shows the effect does not require an ensemble:

\begin{center}\small
\begin{tabular}{cccc}
\toprule
$K$ & monitor MAE $\downarrow$ & $\rho$ $\uparrow$ & ceiling MAE \\
\midrule
1 & 6.908 & 0.533 & 8.897 \\
3 & 6.582 & 0.547 & 7.809 \\
5 & 6.504 & 0.559 & 7.583 \\
\bottomrule
\end{tabular}
\end{center}

At $K=1$ there is no ensemble disagreement at all, only single-model entropy and latent
statistics, and the monitor still beats its ceiling --- by the largest margin of the three.
Ensembling improves absolute accuracy and improves the ceiling faster than it improves the
monitor.

\subsection{Size-matched transfer control}

Table~\ref{tab:transfermatched} repeats the transfer matrix with every regime subsampled to the
A-only training-anchor count ($22{,}200$), so all regimes see equal data.

\begin{table}[h]
\centering\small
\begin{tabular}{lcccc}
\toprule
training regime & test A & test B & test C & test D \\
\midrule
A & \emph{+0.580} & +0.218 & +0.213 & +0.135 \\
A+B & \emph{+0.575} & \emph{+0.747} & +0.278 & +0.759 \\
A+C & \emph{+0.573} & +0.411 & \emph{+0.477} & +0.463 \\
A+D & \emph{+0.575} & +0.694 & +0.287 & \emph{+0.844} \\
A+B+C & \emph{+0.569} & \emph{+0.745} & \emph{+0.466} & +0.762 \\
\bottomrule
\end{tabular}
\caption{Size-matched control for Table~\ref{tab:transfer}. Spearman $\rho$, mean over three
seeds; italic entries are configurations seen in training.}
\label{tab:transfermatched}
\end{table}

The critical cell $\mathrm{A{+}B{+}C}\to$D gives MAE $3.622$ and $\rho = 0.762$, against
$3.584$ and $0.766$ unmatched. The effect is therefore not explained by sample count alone. We
do not claim that diversity is its only cause: severity of the donor shift covaries with its
presence, and we did not manipulate the two independently.

\subsection{Rejected instruments}

Two configurations failed the pre-registered validity gate (a new configuration counts as an
instrument only if $\Hstar$ is at least $40\%$ below A's \emph{and} the no-exposure monitor
transfers to it with $\rho < 0.35$):

\begin{center}\small
\begin{tabular}{llcc}
\toprule
& perturbation & $\Hstar$ & no-exposure $\rho$ \\
\midrule
F  & obs.\ noise $0.15$ & 15.23 & $+0.472$ \\
F2 & obs.\ noise $0.70$ & 12.72 & $+0.359$ \\
F3 & action dropout $0.20$ & \phantom{0}3.67 & $+0.238$ \\
\bottomrule
\end{tabular}
\end{center}

Observation noise never bites: even at $70\%$ cell corruption $\Hstar$ falls only from $15.5$
to $12.7$. An action-conditioned model with sixteen steps of belief context does not depend
heavily on the corrupted channel. F and F2 are reported as invalid instruments rather than as
evidence about transfer; predictions targeting them are uninformative by rule. The severity
ladder used to select F2 and F3 is in \texttt{results/aleatoric/calibration.json}.

A useful control follows from this. The regime $\mathrm{A{+}F2}$ --- exposure to a
configuration that barely degrades the model --- teaches nothing transferable: $\rho = 0.226$
on B and $0.087$ on D, at or below the no-exposure baseline. Exposure helps only when the donor
genuinely degrades the world model.

\subsection{Planning experiment (weak test)}

The adaptive-planning arm sets the imagination horizon of a random-shooting model-predictive
controller to the monitor's prediction, clipped to $[1, 30]$, and is charged for the probe
rollout used to produce that prediction ($K \cdot H_{\max}$ imagined states per decision).
Candidate action sequences are drawn from the same forward-biased distribution used to collect
training data. Targets require a detour: shortest-path distance must exceed Manhattan distance.
Sixty layouts, paired across arms.

\begin{center}\small
\begin{tabular}{lccc}
\toprule
arm & success & imagined states / decision & mean $H$ \\
\midrule
fixed $H{=}3$  & $0.917 \pm 0.036$ & \phantom{0}288 & \phantom{0}3.0 \\
fixed $H{=}5$  & $0.933 \pm 0.032$ & \phantom{0}480 & \phantom{0}5.0 \\
fixed $H{=}10$ & $0.767 \pm 0.055$ & \phantom{0}960 & 10.0 \\
fixed $H{=}20$ & $0.517 \pm 0.065$ & 1920 & 20.0 \\
fixed $H{=}30$ & $0.367 \pm 0.062$ & 2880 & 30.0 \\
adaptive       & $0.517 \pm 0.065$ & 2637 & 25.9 \\
\bottomrule
\end{tabular}
\end{center}

Success falls monotonically with horizon, so the compute--performance frontier has no interior
optimum and an adaptive rule has nothing to find. The detour manipulation did not take: mean
detour was $0.13$ cells, with $57$ of $60$ layouts falling through to the least restrictive
tier. We therefore report this as a weak test of downstream utility rather than a refutation.
The adaptive arm's specific failure --- selecting $H \approx 26$ where $3$--$5$ is optimal ---
is the anti-conservative bias of \S\ref{sec:ood} reappearing under the state distribution that
model-predictive control itself induces.

\subsection{Information-boundary verification}

The monitor may use no observation, state, or error at time $>t$. This is verified
executably rather than asserted: every field of the evaluation window strictly after the anchor
is replaced with random values, the features are recomputed, and all feature columns must be
bit-identical while the target must change. Corrupting the future leaves the features
unchanged and moves mean $\Hstar$ from $6.45$ to $0.01$. The planned action sequence is
deliberately preserved by the corruption, since a planner genuinely knows it.
